\section{Reading the congruence length on the localized column}

% BEGIN E78R1
\paragraph{E78R1. The cut in the rational source before reading the index.}
Fix the real and geometric datum of a provisional construction state,
with its original face, phase-one pins, primitive and auxiliary row
conditions ($\mathrm{Holds}$), reader $Y$, and rational reader $Z$
with its congruence field still omitted. Write $K$ for its coefficient
field. Fix a row $i$, put $j=i+1$, and fix a tuple $t$. Write
$e=Y.eMt(i,t)$ for the rational coordinate isomorphism,
$J=j_{Mt}e$ for its complex reading, $H$ for the coordinate
Hermitian matrix and $E,\Pi$ for the source's cut inclusion and
projection. Write $e_g$ for the analytic good idempotent and
$\Theta$ for the stored theta map. The comparison packet gives
\[
 \Pi E=1,\qquad E^*H(1-E\Pi)=0,\qquad
 \Theta Jx=j_W(M\Pi x),
\]
where $M$ is the injective base lift matrix for this tuple; the images
under $j_W$ of the target basis are complex-linearly independent.
For clarity, the last equality uses no index: the integral reading
compares the lift on $E\Pi x$ with the base matrix. On
$x-E\Pi x$ the primitive orthogonality relation kills the projected
lift. The phase is one and the target evaluation has positive norm
factor, so equality after evaluation gives equality in the target.
These are the two parts of the base theta calculation in BTH7--BTH11
on the provisional reader. The input and target norm, orthogonality,
frame and evaluation identities used there all occur before its
congruence field.

Extend the coordinate projection via the complex basis $J(e_a)$.
Its kernel is the orthogonal complement of the span $C$ of the
columns $J(Ee_k)$, by the first two identities. In the last identity
both sides extend complex-linearly in that basis; independence of
the target columns and injectivity of $M$ after scalar extension
give $\ker\Theta=C^\perp$. The good idempotent fixes every column
of $E$ by its level action. Consequently
\[
 Z_0=(\ker\Theta)^\perp\cap\operatorname{im}e_g=C .
\]
For the rational source $V$, put $U=j_{Mt}^{-1}(Z_0)$ and
$U'=j_{Mt}^{-1}(Z_0^\perp)$. Then $V=U\oplus U'$,
$U=eE(K^{\operatorname{cols}(E)})$ and the rational projection
$p$ onto $U$ sends $ex$ to $eE\Pi x$. Indeed $ex-eE\Pi x$
has reading orthogonal to $C$, and $j_{Mt}$ is injective by
independence of its basis readings. This also shows that a zero cut
or a zero-dimensional row requires no exceptional choice.
% END E78R1

% BEGIN E78R2
\paragraph{E78R2. The two geometric lattices read the rational quotient.}
Let $N$ be the original integral coordinate lattice and $L$ the
$\mathcal O_K$-span of the rational preimages of the specified analytic
lattice generators. The integral reading identifies $eN=L$.
Indeed each specified generator has a rational preimage, and the
membership test puts its coordinates in $N$. Conversely the
reading of $ex$, for $x\in N$, is an integral combination of those
generators. Lifting them and using injectivity of $j_{Mt}$ expresses
$ex$ as that same combination in $L$.

The numerator and denominator lattices defining the rational congruence
index are
\[
 A=U\cap(L+U')=pL,\qquad B=L\cap U.
\]
The first equality with $pL$ follows by decomposing each $l\in L$
along $U\oplus U'$. The injection $q=eE$ therefore sends $\Pi N$
onto $A$. It sends $\{z:Ez\in N\}$ onto $B$, since $e$ is an
isomorphism and $E$ is injective. These two domain lattices are the
source's $T_l$ and $\mathrm{cutZ}$, respectively, by the
cut-lattice identity on this block. Thus $q$ identifies their
quotient with $A/B$. Localization preserves this isomorphism;
in particular the normalized length $\mathrm{etaRat}(D,t,j,v)$ is
exactly the geometric family's $\eta_J$ at this tuple and place.
Normalization here divides the localized length by the length of
the local integer ring modulo the residue characteristic.
The order of the lattices matters: rational orthogonal projection
need not preserve an integral lattice, and the quotient measures
precisely the difference between its projected and intersected lattices.
% END E78R2

% BEGIN E78R3
\paragraph{E78R3. Integral projectors for residual Hecke classes.}
At a twisted place fix the normalized local comparison of the
source and the twisted Hecke hypothesis $\mathrm{HeckeStdTw}$.
Let $O$ be its coefficient DVR with fraction field $F$, and
$\mathbb T$ the $O$-algebra generated by its commuting operators on
the finite-dimensional $F$-space. Their split generalized eigenvalues
lie in $O$. There are only finitely many supported joint characters:
successively split a common generalized eigenspace whenever a
generator has two eigenvalues on it. Each nontrivial split increases
the number of nonzero blocks, bounded by dimension. On a terminal
block every generator is scalar modulo the commuting nilpotent
operators, which gives its integral joint character. This argument
also applies to an infinite generating set.

For a fixed supported residual class $\mathfrak m$, choose for every
supported character $\phi$ outside it an element of its residual
maximal ideal outside $\mathfrak m$. Their product $a\in\mathbb T$
has unit eigenvalue modulo the maximal ideal of $O$ on the chosen
class and zero eigenvalue modulo that ideal on all remaining classes.
Form $g,h\in O[X]$ from the factors $(X-\phi(a))^d$ for the
chosen and remaining characters, respectively, with $d$ large enough
to kill the nilpotent part on each block. Roots from opposite sets
differ by a unit, so $g,h$ are comaximal in $O[X]$. If
$bg+ch=1$, the operator $(ch)(a)$ acts as the identity on the
first blocks and as zero on the others. It lies in the integral
Hecke algebra. For an unsupported class take the zero projector.
The supported class projectors are disjoint and sum to one by the
generalized eigenspace decomposition. If the row lattice is stable
under the generators, these operators split that very lattice into
its residual blocks, even though eigenvalues within one residual
class may coincide after reduction.

Here is one source-level criterion for the stability used above.
Assume $Q( w_i)=Q_i$ and the original $\mathrm{LatInt}$ at the
twisted place: the rational generators lift and every designated
Hecke translate of each generator has coefficients integral at that
place. The local identification of the rational with the analytic
operator and injectivity of the reading give exactly that sum on
its chosen rational preimages. Embed the coefficients in $O$ and
extend by $O$-linearity to the whole local lattice. When there is a
localization above $(j,v)$ its $Q_i$ identifies the exclusion set
with the Hecke model's set. When there is none, the exhaustive
maximal-ideal list forces $\mathbb T$ to be the zero ring, since any
nonzero commutative ring with one has a maximal ideal. Its identity
acts as zero, so the ambient space is zero and stability holds there
as well.
% END E78R3

% BEGIN E78R4
\paragraph{E78R4. Summing the actual local lengths.}
Fix $1\leq j\leq m$ and a twisted place $v$ of residue
characteristic $\ell$. Keep the normalized local comparison, twisted
Hecke condition and lattice stability of E78R3. In addition assume
the source's spectral projection assertion $\mathrm{SpecPr}$ at the
localizations $a$ over $(j,v)$: projecting a rational vector of
$U$ lands in the sum of the good point eigenspaces there, and
projecting a vector of $U'$ lands in the complementary point sum.
Use the comparison embeddings $\Psi_a$ into the common $F$-space.
Denote the two images of these sums by $G_a,B_a$, and the image of
the original local integral module $\mathrm{MtIn}_a$ by $L_a$.

The $\Psi_a$ have range the specified residual blocks. Their
differences from the projection of a rational vector are in the
off-block spaces, so projection by the idempotents of E78R3 identifies
the sum of those projected vectors with the original rational vector
after extension. The local integral comparison in
$\mathrm{HeckeStdTw}$ identifies $L_a$ with the intersection of its block
and the sum of the local lattice and the off-block space. This is the image of
the integral projector on the lattice, by the calculation
$U\cap(L+U')=pL$ in E78R2. Thus the extended lattice is
$\bigoplus_a L_a$. The distinct point eigenspaces inside one block
are independent: for a finite dependence, an operator separating
two characters removes one term after subtraction of its eigenvalue,
and induction proves directness of the character spaces. Distinct
residual blocks are already independent. By $\mathrm{SpecPr}$ the
extensions $U_F,U'_F$ lie in $\sum_a G_a$ and $\sum_a B_a$.
The latter sums are disjoint, and $U_F+U'_F$ is the entire space
by E78R1 and the comparison basis. Both inclusions are therefore
equalities. In particular the two local point sums fill their block.

Since a torsion-free module over a DVR is flat, extension from
$\mathcal O_{K,v}$ to $O$ preserves the intersection with the cut
(as a kernel) and the image of its projection. The comparison basis
identifies the extended ambient vector space with the prescribed
one. Direct sums now give an isomorphism of $O$-modules
\[
 (A/B)_v\otimes_{\mathcal O_{K,v}}O\ \simeq \bigoplus_a \frac{G_a\cap(L_a+B_a)}{L_a\cap G_a}.
\]
The summand is the original localized quotient
$\mathrm{MtUp}_a/\mathrm{MtLow}_a$ via the injective $\Psi_a$ and
its coefficient-ring isomorphism. Each has finite length: $L_a$
spans its block and is finitely generated, its projection is finitely
generated, and a nonzero denominator sends each projected generator
into $L_a\cap G_a$. A finitely generated torsion module over a DVR
has finite length.

Write $e$ for the valuation in $O$ of a uniformizer of
$\mathcal O_{K,v}$. It is positive since this is the specified local
embedding. A composition series over $\mathcal O_{K,v}$ shows that
a finite length multiplies by $e$ upon extension; the same applies
to the length of the ring modulo $\ell$. The localizations' residue
characteristic is the same $\ell$, so their normalized lengths have
one denominator in $O$. Additivity of length and cancellation of $e$
yield
\[
 D.\eta_J(j,v)=\sum_{a\in D.\mathrm{locs}(j,v)}\eta_{\rm loc}(a)
 =\mathrm{etaRat}(D,\mu_{\rm tw}(v),j,v).
\]
This proof retains the integral residual blocks. Dividing the
characteristic-zero eigenspaces further while keeping arbitrary
integral lattices would not justify the displayed direct sum.
% END E78R4

% BEGIN E78R5
\paragraph{E78R5. The eta field and the fitted localized reader.}
Let $x$ be an original provisional state on $P$, and let
$\mathcal L$ be a localization of its coefficient data. On
$D'=\mathrm{withLoc}(\mathcal L)$ assume a normalized local comparison
for the indicated $M$, twisted standard Hecke action, the $Q$ identities
and twisted $\mathrm{LatInt}$ of E78R3, and the $\mathrm{SpecPr}$
of E78R4 at each existing localization. The old datum, reader,
primitive rows and auxiliary rows are those of $x$. We claim the
source's guarded identity at this localization, simultaneously in
every $1\leq j\leq m$ and place $v$.

The tuple reader sends the reference tuple to the reference tuple,
reads the marked tuple at each place, and has a left inverse.
Its injectivity together with the tuple pins therefore identifies a
nonreference tuple with exactly the twisted places. Off this set no
index belongs to $D'.\mathrm{locs}(j,v)$, whose members by definition
are twisted. The defining sum is empty and gives zero, as does the
geometric formula. At a twisted place the tuple is the marked
$\mu_{\rm tw}(v)$. E78R1 uses the phase-one pins and primitive data
from $x$ there before using any eta law. Its lattice reading is
$x.r.Y.\mathrm{latL}$, so E78R2 and E78R4 identify $D'.\eta_J$
with the definite index on the actual moved realization. Writing
$i=j-1$ in $\mathrm{Fin}(m)$ gives exactly
$\mathrm{etaOf}(x.m.z.R,x.m.z.F;j,v)$ in both cases. The all-place
copy identity of PKT29 now gives
\[
 D'.\eta_J(j,v)=P.q.\mathrm{eta}(j,v).
\]
This also proves the version with a transported place $v'$ after
substitution of its dependent equality with $v$.

The later stage prescribed by the source is
$o_{\mathcal L}=(\eta_{\mathcal L},V,S,N_{\mathcal L},G)$, where
$\eta_{\mathcal L}$ is the localization's own fixed congruence function.
Its assembled $D$ equals $D'$ by construction. It satisfies
$\mathrm{Fits}$ for every localization of this state: the division
stage retains its period and universal division laws, the place,
Teichm\"uller, character and series clauses of its group are the old
ones, the place map and central idempotent are unchanged, and the
central actions on ambient and lift spaces are the ones already
read by $x$. Thus all four defining fields are supplied. Use the identity on the finite group of places to move the datum
of $x$ to this stage: compose its packet copy with that identity.
The core and real data transports compose, so this moved datum is
exactly the prescribed $\mathrm{datL}(x,\mathcal L)$. The unchanged
comparison maps give the nine coherences of PKT30. For the reader
construction now assume the original analytic $\mathrm{Inputs}(D')$
and apply RDR12 and RDR20--RDR21 to choose the rational side and
pinned reader there. The guarded eta equality above was proved
before this assumption. It
gives precisely the last field of RDR22. This supplies the full
rational reader and $\mathrm{RowsAre}$ on the localized state, with
the same tuple map and its marked inverse. The construction starts
from $x$ and the specified local comparison. Establishing the
required good-point comparison and supplying a suitable initial state
throughout the selected family remain separate producers.
% END E78R5

% BEGIN E78R6
\paragraph{E78R6. How the good-point comparison supplies the spectral projection.}
At each localization $a$ retain the source's $\mathrm{HeckeTied}$,
$\mathrm{HeckeRational}$ and $\mathrm{AVTied}$ data for $M$, split
uniformizers, spanning by point eigenspaces and the internal isotypic
decomposition. At every tuple and row $(\mu(a),j(a))$ used below,
retain the provisional reading of E78R1, or assume the native
$\mathrm{ThetaCut}$ there, namely
$\operatorname{im}\Theta^*\subseteq\operatorname{im}e_g$.
Assume that the spherical Hecke operators act by scalars
on its automorphic type spaces and that distinct type spaces there
are orthogonal. Use the native $\mathrm{ConclusionZ}$ identity for
the image under $\Theta^*$ of the good rows. Assume also its row
coverage hypothesis: each nonzero target isotypic component at $j$
is good and occurs at a localization of $(j,v)$ satisfying the
good-point guard. Under these assumptions the $\mathrm{SpecPr}$ of
E78R4 holds. Here are the two parts of the proof.

First $(\ker\Theta)^\perp=\operatorname{im}\Theta^*$ by the
finite-dimensional adjoint identity. E78R1 puts this subspace
in $\operatorname{im}e_g$ at the stated tuple, as does the native
cut premise directly. Its intersection with $\operatorname{im}e_g$
is consequently the entire subspace, so $Z_0=\operatorname{im}\Theta^*$.
Transport the automorphic type
spaces at all good points over $(j,v)$ to the chosen localization
$a$. This transport uses equality of their tuples and row indices.
The isotypic function space depends only on the equivalence class
of the irreducible representation; the realization is injective and
the catalogue covers every occurring irreducible type. Thus a
transported nonzero type is an actual member of the catalogue at $a$.
The good-row image identity puts each such type in $Z_0$. Conversely,
decompose the target into isotypic components before applying
$\Theta^*$. For each nonzero component, row coverage selects a
localization, and the same good-row image identity puts its image
in the sum of the transported types. Hence those types span $Z_0$,
and all other types lie in $Z_0^\perp$ by orthogonality. The
catalogue spans the full space by $D.\mathrm{naut}.\mathrm{typeV\_top}(a)$.
Express a vector of $Z_0^\perp$ as a finite sum from these type
spaces and collect the terms from the transported types. Their
sum belongs to $Z_0$ and, after subtracting the remaining terms,
to $Z_0^\perp$. Hence their sum vanishes. This proves that the
remaining type spaces span all of $Z_0^\perp$ as required for the
second projection.

A type with a marked good point belongs to this span using its own
localization. A type marked by a bad point cannot coincide with a
transported good type. At the same localization nonzero independent
type spaces can coincide only for the very same type, whose point
is uniquely specified. At a different localization the Hecke tie
provides a split operator on which the two point eigenvalues differ
by a unit under the prescribed coefficient identification. Equality
of the transported type spaces would force equality of these
eigenvalues: act on a nonzero vector of that space and use
injectivity of its realization and coefficient embedding. Thus the
marked points have the required signs for the two sums, even in the
presence of types without a marked point.

Second, if the complex reading of a rational vector $y$ lies in the
sum of a set of types $\mathcal S$, its local projection can have
components only at points marked by a member of $\mathcal S$.
Suppose it had a component at another point $b$. Since the point
eigenspaces span and are independent, choose a linear functional
on the local vector space vanishing on all other point spaces and
nonzero on this projection. It is a dual Hecke eigenfunctional
with character $b$. Compose it with the local projection and the
coefficient embedding into $\mathbb C$. Extend the resulting
semilinear functional to a complex-linear functional $\lambda$
using the original rational frame as complex basis. Its value on
the reading of $y$ is nonzero, so it is nonzero on some type in
$\mathcal S$. For every allowed spherical operator $T$, the rational
Hecke comparison gives a rational action on frame vectors and
therefore $\lambda T=b(T)\lambda$ in the prescribed embedding.
The analytic operator is the same one there: its finite sum over
cosets is complex-linear and both descriptions agree on the frame.
Evaluate on a vector of that type where $\lambda$ is nonzero. Its
spherical scalar must be $b(T)$ for every allowed $T$. The marking
criterion of $\mathrm{AVTied}$ then marks this type with $b$, a
contradiction. Applying the conclusion to the types in $Z_0$ and
to those in $Z_0^\perp$ gives both parts of $\mathrm{SpecPr}$.

For the lattice assertion take a single coefficient-built witness
$\mathrm{Built}(A,M_0,C,\mathrm{tup},\mathrm{big},L)$ with its
$\mathrm{MtInput}(D,M_0,\mathrm{tup},C.Q)$. On its specialization
$\mathrm{withLoc}(L)$ take $M=\mathrm{mtL}(L,M_0)$ and $Q=C.Q$.
The lattice condition of E78R5 follows for these data. The
specialization identifies its excluded-prime set with the one in
the row. The full-lattice clause supplies rational preimages of all
generators and the integral Hecke clause supplies their required
expansions. This gives the $\mathrm{LatInt}$ of E78R3 at every
twisted $(j,v)$. With the stated geometric, Hecke, and good-point
premises, E78R5 now gives the guarded eta equality on that member.
Its final reader construction also uses $\mathrm{Inputs}(D')$.
% END E78R6
